\documentclass[12pt,a4paper]{article}

\usepackage[utf8]{inputenc}
\usepackage[T1]{fontenc}
\usepackage[brazil]{babel}
\usepackage{geometry}
\usepackage{xcolor}
\usepackage{listings}
\usepackage{hyperref}

\geometry{margin=2.5cm}

% Os números ficam na margem interna do quadro, sem ultrapassar sua largura.
\lstdefinestyle{codigo}{
	numbers=left,
	numberstyle=\tiny\color{gray},
	numbersep=6pt,
	xleftmargin=2.5em,
	frame=single,
	framexleftmargin=2em,
	framesep=5pt,
	breaklines=true,
	breakatwhitespace=false,
	columns=fullflexible,
	keepspaces=true,
	basicstyle=\ttfamily\small,
	captionpos=b
}
\lstset{style=codigo}

\title{Título do relatório}
\author{%
	\begin{tabular}{rl}
		Nome: & \rule{9cm}{0.4pt} \\
		Matrícula: & \rule{9cm}{0.4pt}
	\end{tabular}%
}
\date{\today}

\begin{document}
\maketitle

\section{Introdução}
Apresente aqui o objetivo e o contexto do relatório.

\section{Metodologia}
Descreva os métodos utilizados. Para incluir código com numeração de linhas:

\begin{lstlisting}[language=Python, caption={Exemplo de código}, label={lst:exemplo}]
def soma(a, b):
		return a + b
\end{lstlisting}

Consulte o Código~\ref{lst:exemplo} ao discutir a implementação.
Também é possível incluir um arquivo externo com
\verb|\lstinputlisting[language=Python, caption={...}, label={...}]{caminho/arquivo.py}|.

\section{Resultados e discussão}
Apresente e discuta os resultados.

\section{Conclusão}
Resuma as conclusões do trabalho.

\end{document}
